\subsection{导子与代数同态之间的关系}

我们在本号中再次假设第 2 号的情形 (II) 成立，且不假设分次 K-代数 A 是结合的。给定一个元素 \(\delta \in \Delta\) ，考虑分次 K-模 E( \(\delta\) ) （II，§ 11，第 2 号），它使得

\[
(\mathbf {E} (\delta)) _ {\mu} = \mathbf {E} _ {\mu + \delta}
\]

对所有 \(\mu \in \Delta\) 成立。我们在分次 K-模 \(\mathbf{A} \oplus \mathbf{E}(\delta)\) 上定义一个分次 K-代数结构：对每个齐次元素 \(a \in \mathbf{A}\) 以及任意元素 \(a' \in \mathbf{A}, x, x'\) （ \(\mathbf{E}(\delta)\) 中的），规定

\[
(a, x) \left(a ^ {\prime}, x ^ {\prime}\right) = \left(a a ^ {\prime}, x. a ^ {\prime} + \varepsilon_ {\delta, \deg (a)} a. x ^ {\prime}\right);\tag{28}
\]

验证这个乘法给出一个分次环结构是直接的。

投影 \(p:(a,x)\mapsto a\) 称为代数 \(\mathbf{A}\oplus\mathbf{E}(\delta)\) 的增广（映射），并且是一个分次代数同态。分次 K-线性映射 \(g:\mathbf{A}\to\mathbf{A}\oplus\mathbf{E}(\delta)\) （次数为 0 ）使得复合

\[
\mathrm{A} \xrightarrow {g} \mathrm{A} \oplus \mathrm{E} (\delta) \xrightarrow {p} \mathrm{A}
\]

是恒同映射 \(l_{A}\) 的那些映射，正是形如 \(x \mapsto (x, f(x))\) 的映射，其中 \(f: A \to E\) 是次数为 \(\delta\) 的分次 K-线性映射。

命题 12. 为使分次 K-线性映射 \(f:A\to E\) （次数为 \(\delta\) ）是一个 \(\varepsilon\) -导子，必要且充分条件是：映射 \(x\mapsto(x,f(x))\) （从 A 到 \(\mathbf{A}\oplus\mathbf{E}(\delta)\) ）是一个分次 K-代数同态。

利用如下事实：对齐次的 \(x\) （ \(\mathbf{A}\) 中的）以及 \(y \in \mathbf{A}\) ，

\[
(x y, f (x y)) = (x, f (x)). (y, f (y)),
\]

并考虑到(28)，我们得到等价关系

\[
f (x y) = f (x). y + \varepsilon_ {\delta , \deg (x)} x. f (y),
\]

由此得出该命题。

命题 13. 为使代数 \(\mathbf{A} \oplus \mathbf{E}(\delta)\) 是结合的且有单位元的，必要且充分的是： \(\mathbf{A}\) 是结合的且有单位元的，并且映射 \((a, x) \mapsto a.x\) 和 \((a, x) \mapsto x\) . \(a\) 在 \(\mathbf{E}\) 上定义一个 \((\mathbf{A}, \mathbf{A})\) -双模结构； \(\mathbf{A} \oplus \mathbf{E}(\delta)\) 的单位元则为 \((1, 0)\) 。

若把元素 \((u,m)\in\mathbf{A}\oplus\mathbf{E}(\delta)\) 写成该代数的单位元，则立即发现 u 必为 A 的单位元；写出 \((u,m).(0,x)=(0,x).(u,m)=(0,x)\) ，我们得到 u.x=x.u=x （对 \(x\in E\) ），并且写出 \((u,m).(u,0)=(u,0).(u,m)=(u,0)\) ，我们得到 m=0。当 \(\mathbf{A}\oplus\mathbf{E}(\delta)\) 结合时，A 的结合性由增广是一个满同态这一事实得出。条件 \((x.a') .a''=x.(a'.a'')\) 等价于 \(((0,x)(a',0))(a'',0)=(0,x)((a',0)(a'',0))\) ，类似地，条件 \(a.(a'.x)=(a.a').x\) 等价于

\[
(a, 0) ((a ^ {\prime}, 0) (0, x)) = ((a, 0) (a ^ {\prime}, 0)) (0, x);
\]

最后，条件 \(a \cdot (x \cdot a') = (a \cdot x) \cdot a'\) 等价于

\[
(a, 0) ((0, x) (a ^ {\prime}, 0)) = ((a, 0) (0, x)) (a ^ {\prime}, 0).
\]

